\documentclass[a4paper,11pt]{ctexart}
\usepackage[left=1.9cm,right=1.9cm,top=2.0cm,bottom=2.0cm]{geometry}
\usepackage{amsmath,amssymb}
\usepackage{enumitem}
\usepackage{fancyhdr}
\usepackage{lastpage}
\usepackage{listings}
\usepackage{xcolor}
\usepackage{array}

\pagestyle{fancy}
\fancyhf{}
\renewcommand{\headrulewidth}{0pt}
\fancyfoot[C]{第~\thepage~页\quad 共~\pageref{LastPage}~页}

\lstset{
  language=C,
  basicstyle=\ttfamily\small,
  columns=fullflexible,
  keepspaces=true,
  breaklines=true,
  showstringspaces=false,
  tabsize=4,
  aboveskip=3pt,
  belowskip=3pt,
  xleftmargin=1.5em,
  escapeinside={(*@}{@*)},
  commentstyle=\color{black!70}
}

\newcommand{\blankline}[1]{\underline{\hspace{#1}}}
\newcommand{\ansbox}{\hfill(\qquad\quad)}
\newcommand{\secbox}[2]{%
  \noindent\begin{tabular}{|>{\centering\arraybackslash}m{1.5cm}|>{\centering\arraybackslash}m{1.5cm}|}
    \hline
    分数 & \\ \hline
    评卷人 & \\ \hline
  \end{tabular}\hspace{0.8em}%
  \begin{minipage}[c]{\dimexpr\textwidth-4.2cm\relax}
    \textbf{#1}#2
  \end{minipage}\par\vspace{0.5em}%
}

\begin{document}

\begin{center}
  {\large\textbf{华中科技大学人工智能与自动化学院 2022～2023 学年第一学期}}\\[4pt]
  {\Large\textbf{“C 语言程序设计”考试试卷（A 卷）}}\\[8pt]
  考试方式：\underline{\quad 闭卷 \quad}\qquad
  考试日期：\underline{\quad 2023-02-14 \quad}\qquad
  考试时长：\underline{\quad 150 \quad} 分钟
\end{center}

\vspace{4pt}
\begin{center}
  \begin{tabular}{|>{\centering\arraybackslash}p{1.6cm}|*{7}{>{\centering\arraybackslash}p{1.5cm}|}}
    \hline
    题号 & 一 & 二 & 三 & 四 & 五 & 六 & 总分 \\ \hline
    得分 & & & & & & & \\[8pt] \hline
  \end{tabular}
\end{center}

\begin{center}
  \small 院系：\underline{\hspace{2.3cm}}\quad 专业班级：\underline{\hspace{2.2cm}}\quad 姓名：\underline{\hspace{2.0cm}}\quad 学号：\underline{\hspace{2.3cm}}
\end{center}
\vspace{6pt}

% ==================== 一 ====================
\secbox{一、判断下列语句或程序段的对错。}{（“$\times$”表示错，“$\checkmark$”表示对）（10 分）}

\begin{enumerate}[label=(\arabic*), leftmargin=2.4em, itemsep=3pt]
  \item \texttt{int x[10], *px;}\\
        \texttt{px= ++x;} \ansbox
  \item \texttt{unsigned char \_x\_y\_z\_1 = 056;} \ansbox
  \item \texttt{int data[4]=\{0, 1, 2, 3, 4\} ;} \ansbox
  \item \texttt{int a[5][5], (*p)[5];}\\
        \texttt{p= a+2;} \ansbox
  \item \texttt{\#define\quad N\quad 50}\\
        \texttt{int\quad a[N/3+1];} \ansbox
  \item \texttt{int x ='A' +'B';} \ansbox
  \item \texttt{char str[15]= "Hello, HUST! ", *p="Hello, AIA! ";}\\
        \texttt{strcpy(p,str);} \ansbox
  \item \texttt{register int a;}\\
        \texttt{scanf("\%d", \&a);} \ansbox
  \item \texttt{int a[5], (*p)[5];}\\
        \texttt{p=(int (*)[5])a;}\\
        \texttt{p[0] = 10;} \ansbox
  \item \texttt{char *str = "\textbackslash\textbackslash\textbackslash"ABC\textbackslash"\textbackslash\textbackslash\textbackslash";} \ansbox
\end{enumerate}

\vspace{8pt}

% ==================== 二 ====================
\secbox{二、计算下列表达式的值}{（10 分）}

\noindent 假设 \texttt{unsigned int a=2, b=3; int c=4, d=5; float k=3.0, f;}\\
计算下列各表达式的值，假设 \texttt{int} 类型为 16 位长度，且各题彼此独立

\begin{enumerate}[label=(\arabic*), leftmargin=2.4em, itemsep=3pt]
  \item \texttt{f=(d<=-c+2\&\&!c+1.0)/2;} \ansbox
  \item \texttt{k= k/a + c/d;} \ansbox
  \item \texttt{d= (a-b)>0 ? (a-b) : -(a-b);} \ansbox
  \item \texttt{a += b \%= a+b;} \ansbox
  \item \texttt{a\^{}(b\^{}a)\&(a|d);} \ansbox
\end{enumerate}

\vspace{8pt}

% ==================== 三 ====================
\secbox{三、程序改错（10 分）}{要求：不得改变程序框架，不得重写程序，无需文字说明，直接在代码上添加、删除和修改。}

\noindent (1) 对生成给定大小的二维数组进行转置，数组元素为 0～100 的随机数。
\begin{lstlisting}
#include<stdio.h>
#include<time.h>
#include<stdlib.h>

int main()
{
    int a[N][N], i, j;

    srand((unsigned) time(NULL));
    printf("array a:\n");
    for (i = 0; i < N; i++)
    {
        for (j = 0; j < N; j++)
        {
            a[i][j] = rand() ;
            printf("%4d", a[i][j]);
        }
        printf("\n");
    }
    trans(a);
    printf("transposed array a:\n");
    for (i = 0; i < N; i++)
    {
        for (j = 0; j < N; j++)  printf("%4d", a[i][j]);
        printf("\n");
    }
    return 0;
}

void trans(int **a)
{
    int t, i, j;
    for (i = 0; i < N; i++)
        for (j = 0; j < N; j++)
        {
            t = a[i][j];
            a[i][j] = a[j][i];
            a[j][i] = t;
        }
}
\end{lstlisting}

\noindent (2) 求方程 $ax^2+bx+c=0$ 的根（5 分）
\begin{lstlisting}
#include<stdio.h>

double f1(double a, double b, double Disc, double x1)
{
    double x2;
    x1 = (-b + sqrt(Disc)) / (2 * a);
    x2 = (-b - sqrt(Disc)) / (2 * a);

    return x2;
}

double f2(double a, double b)
{
    return ( (-b) / (2 * a) );
}

void main()
{
    double a,b,c;
    double x1,x2;
    printf("please input a,b,c\n");
    scanf("%f%f%f", &a, &b, &c);

    Disc = b*b - 4 * a*c;
    if (Disc > 0)
    {
        x2 = f1(a, b, Disc, x1);
        printf("x1=%f\nx2=%f\n", x1, x2);
    }
    else (Disc==0)
    {
        x1 = f2(a, b);
        printf("x1=x2=%f\n", x1);
    }
    else
    {
        printf("No real roots\n");
    }
}
\end{lstlisting}

\vspace{8pt}

% ==================== 四 ====================
\secbox{四、程序填空}{（10 分）}

\noindent (1) 编程实现在字符串 \texttt{s1} 中的指定位置 \texttt{n} 处插入字符串 \texttt{s2}。
\begin{lstlisting}
#include <stdio.h>
#include <malloc.h>
(*@\blankline{6.5cm}@*)
{
    int i = 0;
    for (i = 0; i < n; i++)
        (*@\blankline{4.5cm}@*)
    while ( (*@\blankline{4.5cm}@*) )
        *ss++ = *s2++;
    while ( (*@\blankline{4.5cm}@*) )
    {
        *ss++ = *s1++;
    }
    (*@\blankline{4.5cm}@*)
    return 0;
}
int main()
{
    char s1[] = "123456789";
    char s2[] = "1234";
    (*@\blankline{4.5cm}@*)
    insert(s1, s2, s3, 4);
    printf("%s\n", s3);

    return 0;
}
\end{lstlisting}

\noindent (2) 找出一个 $2\times 3$ 的矩阵中的最大值及其行下标和列下标。
\begin{lstlisting}
#include <stdio.h>
int FindMax(int *p, int m, int n, int *pCol);
int main()
{
    int a[2][3], i, j, row, col, max;
    int *p;
    printf("Input 6 datas:");
    for( (*@\blankline{5.5cm}@*) ; p++)  scanf("%d", p);
    printf("\n****array****\n");
    for (i = 0; i < 2; i++)
    {
        for (j = 0; j < 3; j++)
            printf("%4d", a[i][j]);
        printf("\n");
    }
    (*@\blankline{6.0cm}@*)
    printf("\nMax=%d, Row=%d, Col=%d\n", a[row][col], row, col);
}
int FindMax(int *p, int m, int n, int *pCol)
{
    int max, row=0, col=0, i, j;
    (*@\blankline{4.5cm}@*)
    for (i = 0; i < m; i++)
        for (j = 0; (*@\blankline{4.5cm}@*) )
        {
            if (max < *p)
            {
                max = *p;
                row = i;
                col = j;
            }
        }
    (*@\blankline{4.5cm}@*)
    return row;
}
\end{lstlisting}

\vspace{8pt}

% ==================== 五 ====================
\secbox{五、输出程序运行结果}{（25 分）}

\noindent (1)
\begin{lstlisting}
#include <stdio.h>
int f();
int n=1;
void main()
{
    int i,a=0;

    for(i=0;i<5;i++)
        a+=f();
    printf("sum:%d\n",a);
}
int f()
{
    static int i;
    int s=0;
    n *= (i++ +1);
    s+=n;
    printf("i:%d--s:%d\n",i,s);
    return s;
}
\end{lstlisting}

\newpage
\noindent (2)
\begin{lstlisting}
#include <stdio.h>
#include <string.h>
char * pro_str(char *s);
void main()
{
    char string[80];
    printf("input string:\n");
    gets(string);
    puts(pro_str(string));
}
char *pro_str(char *s)
{
    char *temp = s;
    for(;*s!='\0';)
    {
        if(!(*s>='0'&&*s<='9'))
            strcpy(s,s+1);
        else
            s++;
    }
    return temp;
}
\end{lstlisting}
\noindent 输入：\texttt{There are 20 boys and 7 girls in this room.}

\enlargethispage{1\baselineskip}
\vspace{4pt}
\noindent (3)
\begin{lstlisting}
#include <stdio.h>
struct Key
{
    char *keyword ;
    int keyno ;
};

void main( )
{
    struct Key  kd[3] = {{"are", 123}, {"your", 456}, {"my", 789}};
    struct Key  *p =kd;
    int a ;
    char *str ;

    str = (++p)->keyword ;
    printf("str=%s\n", ++str) ;
    a = ++p++ ->keyno %10;
    printf("a=%d", a) ;
    a=p->keyno/100;
    printf("%d\n", a) ;
}
\end{lstlisting}

\vspace{4pt}
\noindent (4)
\begin{lstlisting}
#include <stdio.h>
void main()
{
    char nn[4][5]={"1234","3456","5678","7890"},*pn[4];
    int k,s=0;
    for (k=0;k<4;k++)
        pn[k]=nn[k]+k;
    for (k=0;k<4;k++)
        s=s*10+pn[k][0]-'0';
    printf("%d\n",s);
}
\end{lstlisting}

\vspace{4pt}
\noindent (5)
\begin{lstlisting}
#include <stdio.h>
#include <string.h>
void sortstr(char (*v)[20], int n);
void main( )
{
    char proname[5][20]={"pascal","basic","cobol","prolog","fortran"};
    int i,n;

    n = 'A'-'a';
    sortstr(proname,5);
    for(i=0;i<5;i++)
    {
        *(*(proname+4-i)) += n;
        printf("%s\n",*(proname+4-i));
    }
}
void sortstr(char (*v)[20], int n)
{
    int i,j;
    char temp[20];
    for(i=0;i<n-1;i++)
        for(j=i+1;j<n;j++)
        {
            if(strcmp(v[i],v[j])<0)
            {
                strcpy(temp,v[i]);
                strcpy(v[i],v[j]);
                strcpy(v[j],temp);
            }
        }
}
\end{lstlisting}

\vspace{8pt}

% ==================== 六 ====================
\secbox{六、编写程序（35 分）}{注意：不得使用全局变量，注意程序结构}

\begin{enumerate}[label=(\arabic*), leftmargin=2.2em, itemsep=5pt]
  \item 编程对所输入的分数进行约分，例如，输入 \texttt{72/192}，输出结果为 \texttt{3/8}，若输入 \texttt{72/24}，则输出 \texttt{3}。（7 分）
  \item 已知 $x_1=1$，$x_2=2$，$x_n=0.7x_{n-1}+0.3x_{n-2}$，求 $\{x_n\}$ 前 10 项的平均值（8 分）
  \item 输入一行包含若干单词的字符串，单词之间用空格分开，统计字符串中以“\texttt{ing}”结尾的单词个数，要求功能模块设计合理（10 分）
  \item 建立一个学生班的成绩登记表，包括的信息有班号（全班一个），总人数（设为 20 人），制表日期（整个表一个）和所有学生的信息。每个学生包含下列信息：学号、姓名、性别、四门课程的成绩。要求：(a) 输入登记表及每个学生的上述信息，并统计每个学生的平均成绩；(b) 按平均成绩从高到低的顺序输出排名前 30\% 的每个学生的学号、姓名、平均成绩一览表；(c) 统计平均成绩在 60 分以下学生的人数以及其中男女生的比例，并输出统计结果。编写函数完成上述功能要求，主函数中完成功能函数的调用，平均成绩保留两位小数。（10 分）
\end{enumerate}

\end{document}
